\begin{abstract}
Clinical prediction studies often combine learners into ensembles and explain them with SHapley Additive exPlanations (SHAP), but rarely test whether these explanations are as faithful as those of the single model the ensemble would replace. We compared a hybrid stacked-voting ensemble with the single model chosen by internal cross-validation on eight public healthcare datasets, using leakage-safe repeated cross-validation, pre-written protocols and a deletion-based faithfulness score computed with two feature-removal methods. Against a pre-specified $\pm$0.02 margin, the ensemble's AUC was equivalent overall (pooled difference \Nauc) and on five datasets, and inconclusive on three. Under marginal feature removal its explanations were less faithful on seven of eight datasets, significantly on three. The pooled decrease (\Nndf) was significant with DerSimonian--Laird pooling but not with the more conservative Hartung--Knapp method (95\% CI \Nndfhk), so it is not definitive. Under conditional removal the difference was smaller and, for untuned models, not significant. Similar patterns held on held-out repetitions, with a second SHAP estimator, with tuned models and without excluding patients. XGBoost alone was explained as faithfully as logistic regression; the ensemble was not. Ensemble explanations should therefore be evaluated, with the removal method stated, rather than assumed faithful.
\end{abstract}
